% ==================================================================
% id:       seq-006
% topic:    סדרות
% subtopic: אריתמטיקה של גבולות — הוכח או הפרך
% solution: full
% ==================================================================
\textbf{תהיינה
$(a_n),(b_n)$
סדרות.
    הוכיחו או הפריכו כל אחת מהטענות הבאות:}
\bE
     \item
    אם
    $(a_n)$
    לא מתכנסת
    ואם
    $(b_n)$
    לא מתכנסת, אז
    $(a_n+b_n)$
    לא מתכנסת.
    
\hint{בדקו אם ההתנהגות הלא-מתכנסת של שתי סדרות יכולה להתבטל בסכום. אם נראה לכם שהטענה לא נכונה, חפשו דוגמה נגדית פשוטה.}
\begin{solbox}
    לא נכון. דוגמה נגדית תהיה
    $a_n=(-1)^n, b_n=(-1)^{n+1}$
    כי שתי אלה סדרות שאינן מתכנסות אבל
    $a_n+b_n=(-1)^n-(-1)^n=0$
    זו סדרה קבועה ולכן מתכנסת.
\end{solbox}
    \item
    אם
    $(a_n)$
    מתכנסת
    ואם
    $(b_n)$
    לא מתכנסת, אז
    $(a_n+b_n)$
    לא מתכנסת.
    
\hint{חשבו מה הייתם יכולים להסיק על $b_n$ אילו $a_n+b_n$ הייתה מתכנסת, בעזרת אריתמטיקה של גבולות. אם נראה לכם שהטענה לא נכונה, חפשו דוגמה נגדית פשוטה.}
\begin{solbox}
    נכון. זה נובע מאריתמטיקה של חיסור סדרות מתכנסות!

    נניח בשלילה כי
    $(a_n+b_n)$
    כן מתכנסת, ולכן כחיסור של שתי סדרות מתכנסות נקבל גם כי הסדרה
    $$b_n=(a_n+b_n)-a_n$$
    מתכנסת, בסתירה לנתון שהיא לא.
\end{solbox}
    \item
    אם
    $(a_n+b_n)$
     מתכנסת אז לפחות אחת מבין
    $(a_n)$
    או
    $(b_n)$
    מתכנסות.
    
\hint{מהי השלילה של ״לפחות אחת מהסדרות מתכנסת״?}
\begin{solbox}
לא נכון. זו בדיוק אותה הדוגמה מסעיף (א).
\end{solbox}
      \item
    אם
    $(a_n+b_n)$
    לא
     מתכנסת אז לפחות אחת מבין
    $(a_n)$
    או
    $(b_n)$
    לא
    מתכנסות.
    
\hint{נסחו את הטענה השקולה לה (קונטרה-פוזיטיב) ובדקו אותה בעזרת אריתמטיקה של גבולות.}
\begin{solbox}
    נכון. נניח בשלילה כי שתי הסדרות
    $(a_n),(b_n)$
    מתכנסות, אז מאריתמטיקה של חיבור נקבל כי
    $(a_n+b_n)$
    מתכנסת גם כן, בסתירה להנחה.
\end{solbox}
    \item
    אם
    $a_n<1$
    לכל $n$ ו-$a_n\to L$, אז
    $L<1$.
    \begin{solbox}
לא נכון. דוגמה נגדית:
$a_n=1-\frac1n$.
לכל $n$ מתקיים
$a_n<1$,
אבל
$a_n\to1$,
ו-$1<1$ אינו נכון.
\end{solbox}
    \item
    אם
    $a_n\le1$
    לכל $n$ ו-$a_n\to L$, אז
    $L\le1$.
    \begin{solbox}
נכון. נניח בשלילה ש-$L>1$, וניקח
$\varepsilon=L-1>0$.
לפי הגדרת הגבול, קיים $N$ כך שלכל $n>N$ מתקיים
$|a_n-L|<\varepsilon$,
ובפרט
$a_n>L-\varepsilon=1$,
בסתירה לנתון.
\end{solbox}
    \item
    אם
    $a_n\to L\ne0$
    ו-$(b_n)$
    אינה מתכנסת, אז
    $(a_nb_n)$
    אינה מתכנסת.
\hint{הניחו בשלילה ש-$a_nb_n\to M$, וכתבו את $b_n$ כמנה. כדי לחלק ב-$a_n$ צריך לדעת ש-$a_n\ne0$, לפחות החל ממקום מסוים.}
\begin{solbox}
נכון. נניח בשלילה ש-$a_nb_n\to M$.

\textbf{החל ממקום מסוים $a_n\ne0$:}
נבחר בהגדרת הגבול
$\varepsilon=\frac{|L|}{2}>0$.
קיים
$N$
כך שלכל
$n>N$
מתקיים
$|a_n-L|<\frac{|L|}2$,
ולכן, לפי אי-שוויון המשולש,
$$|a_n|\ge|L|-|a_n-L|>|L|-\frac{|L|}2=\frac{|L|}2>0 .$$

\textbf{סתירה:}
לכל
$n>N$
מתקיים
$b_n=\dfrac{a_nb_n}{a_n}$.
לפי אריתמטיקה של גבולות (מנה של סדרות מתכנסות, כשגבול המכנה
$L\ne0$),
$$\frac{a_nb_n}{a_n}\goes\frac ML .$$
הסדרה
$(b_n)$
שווה לסדרה הזו החל מ-$N$, ולכן, לפי סעיף א' של שאלה
\qref{seq-030},
גם
$b_n\to\frac ML$.
זו סתירה לכך ש-$(b_n)$
אינה מתכנסת.
\end{solbox}
\eE
